
\chapter*{\centering Preface: To My Daughter, on the Eve of Her Gaokao}\addcontentsline{toc}{chapter}{Preface: To My Daughter, on the Eve of Her Gaokao}

One day, I found my daughter reading \emph{Being Digital}.

That was five years ago. In that moment, my heart was torn between joy and sorrow. Joy, because she was finally taking a hard look at the digital world she had grown up in. Sorrow, because when she reached for a book to make sense of that world, the one placed in her hands was still Negroponte's volume from all those years ago.

\vspace{1em}


\emph{Being Digital} prophesied a great migration---from atoms to bits.
Thirty years on, that migration is long complete. Photographs, letters, ledgers, music, transactions, relationships---more and more of what once clung to the physical world has been translated into bits.

\vspace{1em}

But the world did not stop there. Bits can be transmitted, copied, exchanged, encrypted. They become models. They become intelligent.
And so the world took one more step forward: it is no longer just atoms turning into bits---now thought itself can be computed.
This is not some brand-new world. We have been living in it for thirty years. But it is growing more intelligent by the day, and only now is its fuller shape beginning to emerge. I call it \textbf{the Fifth Dimension}.

\vspace{1em}

As someone who works in AI, I witness technology's breakneck advance every day. I also hear, constantly, how hard it is for new graduates to find jobs. A few questions never leave my mind:

The intelligence we are building---\textbf{what does it mean for the young people about to step into the world?}

How do I tell my own child what is still worth holding on to, and what is losing its value?

As ordinary people, what should we be doing to prepare?

\vspace{1em}

These questions set me writing. They have no simple answers.

Optimists say AI will create new opportunities. Pessimists say mass unemployment is inevitable.

But perhaps, before we reach any conclusion, we need to understand what is actually happening.

\vspace{1em}

That is exactly what this book tries to do.
In language that I hope is vivid and fun, it describes six abilities inside the Fifth Dimension that feel almost like magic. It also tells the story of how bits flow; how intelligence takes shape; how the rules of the world are changing; what we should hand over to machines and what we should keep for ourselves; and how, in this new era, we can survive---and become irreplaceable.

\vspace{1em}

My daughter was this book's first reader. She is a high school senior on the humanities track.

She hates probability, cannot keep expected values straight, yet she told me she finally understood what ``entropy'' means---a concept that every information-theory book mentions but that she had never been able to grasp.

She said I ``made technical topics way less stuffy.'' She also said she liked the ``effect before cause'' example---``you know, where you get killed first and \emph{then} see that you're about to get killed''---so cool.

If a math-averse humanities student can follow along, perhaps this book can help a few more people---those students and parents standing on the eve of the storm, feeling a little lost and a little afraid.

\begin{myquote}
This book is for you---\textbf{my daughter, and every ordinary person who, like you and me, refuses to surrender}.
\end{myquote}

In this era of seismic change, there is no magic elixir. There are only towering waves.

The age of AI demands that we stay \textbf{curious}---the way she kept asking ``but what \emph{is} entropy?''

The age of AI demands that we keep our \textbf{fighting spirit}---the way she wrote in her feedback, ``it's not like I want to work in the industry anyway,'' and yet read every page with relish.

\vspace{1em}

May this book be a starting point, and help you find your own Platform Nine and Three-Quarters into the Fifth Dimension.

\textbf{Forward---the answers will come!}

\vspace{3em}

\hfill \myauthorzh

\hfill On the eve of the storm
